
\chapter{Internal Security}

\subsection{Left Wing Extremism (LWE) — From Red Corridor to Naxal-Free Bharat (2014–2025)}

\begin{itemize}

  \item \textbf{Context:} Govt's multi-dimensional counter-LWE strategy targeting complete eradication by \textbf{31 March 2026}; strategy framework: Dialogue $\rightarrow$ Security $\rightarrow$ Coordination.

  \item \textbf{Red Corridor States:} Chhattisgarh, Jharkhand, Odisha, Maharashtra, Kerala, West Bengal, MP, parts of AP \& Telangana.

\end{itemize}

\textbf{Quantitative Decline (2004–14 vs. 2014–24):}
\begin{itemize}
  \item Violent incidents: $\downarrow$53\% (16,463 $\to$ 7,744)
  \item SF deaths: $\downarrow$73\% (1,851 $\to$ 509)
  \item Civilian deaths: $\downarrow$70\% (4,766 $\to$ 1,495)
  \item LWE-affected districts: 126 (2014) $\to$ \textbf{11 (2025)}
  \item Most-affected districts: 36 (2014) $\to$ \textbf{3 (2025)}
  \item Police stations recording Naxal incidents: 330 (76 districts, 2013) $\to$ 52 (22 districts, June 2025)
\end{itemize}

\textbf{Operational Achievements (2024–25):}
\begin{itemize}
  \item \textbf{Operation Black Forest:} 27 hardcore Naxals killed
  \item PLGA battalion abandoned core area in Bijapur-Sukma
  \item Liberated areas: Budha Pahad, Parasnath, Baramasia, Chakrabandha (Bihar); Abujhmad (Chhattisgarh) reached — all after 3 decades of Maoist control
  \item Ops: Octopus, Double Bull, Chakbandha
\end{itemize}

\textbf{Security Infrastructure (Since 2014):}
\begin{itemize}
  \item Fortified Police Stations (FPS): 66 (till 2014) $\to$ \textbf{586} constructed
  \item New security camps: 361 (last 6 years)
  \item Night-landing helipads: 68 constructed
\end{itemize}

\textbf{Financial Choking of Naxals:}
\begin{itemize}
  \item NIA dedicated anti-LWE vertical: 108 cases investigated, 87 charge sheets filed
  \item Assets seized: NIA ₹40 cr + States ₹40 cr + ED attachment ₹12 cr = \textbf{₹92 cr total}
\end{itemize}

\textbf{Key Schemes \& Financial Support:}
\begin{itemize}
  \item \textbf{SRE (Security Related Expenditure):} ₹3,331 cr released to LWE states (11 yrs); 155\% increase over last decade
  \item \textbf{SIS (Special Infrastructure Scheme):} ₹371 cr (SF/SIB); ₹620 cr (246 FPS original phase); extended till 2026 — ₹610 cr (SF/SIB/district police) + ₹140 cr (56 additional FPS)
  \item \textbf{SCA (Special Central Assistance):} ₹3,817.59 cr
  \item \textbf{ACALWEMS:} ₹125.53 cr (camp infrastructure) + ₹12.56 cr (hospitals)
\end{itemize}

\textbf{Infrastructure Development:}
\begin{itemize}
  \item Roads: 12,000 km constructed (May 2014–Aug 2025); 17,589 km approved @ ₹20,815 cr
  \item Mobile towers (Phase 1): 2,343 towers (2G) @ ₹4,080 cr
  \item Mobile towers (Phase 2): 2,542 towers sanctioned @ ₹2,210 cr; 1,154 installed
  \item 4G towers (Aspirational Districts + 4G Saturation): 8,527 approved; 2,596 + 2,761 functional
\end{itemize}

\textbf{Financial Inclusion (90 LWE districts):}
\begin{itemize}
  \item 1,804 bank branches, 1,321 ATMs, 37,850 banking correspondents
  \item 5,899 post offices (every 5 km coverage)
\end{itemize}

\textbf{Skill Development (48 LWE districts):}
\begin{itemize}
  \item 48 ITIs sanctioned @ ₹495 cr; 46 functional
  \item 61 Skill Development Centres (SDCs) approved; 49 functional
\end{itemize}

\textbf{Surrender \& Rehabilitation Policy:}
\begin{itemize}
  \item High-rank cadre: ₹5 lakh; Middle/lower: ₹2.5 lakh; Monthly stipend: ₹10,000 (36 months professional training)
  \item 2025: 521 surrendered; total 1,053 post new state govt formation
\end{itemize}

\textbf{Bastariya Battalion:}
\begin{itemize}
  \item Raised 2018; 1,143 recruits; 400 local youth from Bijapur, Sukma, Dantewada (Chhattisgarh)
  \item Converts former Naxal strongholds into security manpower sources
\end{itemize}
